\documentclass[a4paper, serif, 10pt]{moderncv}

%% moderncv
% style and color
\moderncvstyle{banking}
\moderncvcolor{black}

%% page layout/margins
\usepackage[top=2.5cm, bottom=2.5cm, left=1.5cm, right=1.5cm]{geometry}

\nopagenumbers{}

%% Babel config for French language
% ref:
% - https://latex3.github.io/babel/guides/locale-french.html
\usepackage[french]{babel}

%% fontspec
\usepackage{fontspec}

\IfFontExistsTF{Linux Libertine}{
  \setmainfont[Ligatures={TeX, Common}, Numbers=OldStyle]{Linux Libertine}
}{}
\IfFontExistsTF{Linux Libertine O}{
  \setmainfont[Ligatures={TeX, Common}, Numbers=OldStyle]{Linux Libertine O}
}{}

%% CTeX
% CJK support, used to show CJK characters in the resume
%
% - fontset=none: disable builtin fontset but instead set the CJK font manually
% - heading=false: disable ctex heading
% - punct=kaiming: use kaiming punctuations styles for CJK
% - scheme=plain: use plain scheme, do not override \normalsize font size
% - space=auto: space settings for CJK characters
%
% ref:
% - http://ctan.mirrorcatalogs.com/language/chinese/ctex/ctex.pdf
\usepackage[UTF8, heading=false, punct=kaiming, scheme=plain, space=auto]{ctex}

\IfFontExistsTF{Noto Serif CJK SC}{
  \setCJKmainfont{Noto Serif CJK SC}
}{}
\IfFontExistsTF{Noto Sans CJK SC}{
  \setCJKsansfont{Noto Sans CJK SC}
}{}

%% line spacing
\usepackage{setspace}
\setstretch{1.125}

%% URL styling - use normal text instead of monospace
\urlstyle{same}

% Auto-underline all links
% Defer until \begin{document} so hyperref is already loaded
\AtBeginDocument{%
  \let\oldhref\href
  \renewcommand{\href}[2]{\oldhref{#1}{\underline{#2}}}%
}

%% Patch \httplink and \httpslink to handle full URLs
% The original moderncv macros always prepend http:// or https:// to URLs.
% This causes issues when the URL already has a protocol (e.g., https://example.com)
% resulting in malformed URLs like https://https://example.com.
% This patch checks if the URL already contains "://" and uses it directly if so.
% Note: We use \str_set:Nx (x-expansion) to expand macros like \@homepage first.
\ExplSyntaxOn
\str_new:N \g_yamlresume_url_str

\RenewDocumentCommand{\httpslink}{O{}m}{
  \str_set:Nx \g_yamlresume_url_str {#2}
  \str_if_in:NnTF \g_yamlresume_url_str {://}
    { \url{#2} }
    { \url{https://#2} }
}

\RenewDocumentCommand{\httplink}{O{}m}{
  \str_set:Nx \g_yamlresume_url_str {#2}
  \str_if_in:NnTF \g_yamlresume_url_str {://}
    { \url{#2} }
    { \url{http://#2} }
}
\ExplSyntaxOff

\name{Camille Rousseau}{}
\title{Product Manager Senior — SaaS \& Fintech}
\phone[mobile]{+33 6 12 34 56 78}
\email{camille.rousseau@protonmail.com}
\homepage{https://www.camillerousseau.fr}

\address{12 rue Oberkampf, Paris, Île-de-France, France, 75011}{}{}

\extrainfo{{\small \faLinkedin}\ \href{https://www.linkedin.com/in/camille-rousseau}{@camille-rousseau} {} {} {} • {} {} {}
{\small \faMedium}\ \href{https://medium.com/@camille.rousseau}{@camille.rousseau}}

\begin{document}

\maketitle

\section{Informations de base}

\cvline{}{Product Manager expérimentée avec \textbf{plus de 8 ans d'expérience} dans le développement de produits SaaS et fintech, de l'idéation au lancement à grande échelle.
%
\begin{itemize}
\item Vision produit, feuille de route et priorisation fondées sur la data et les besoins utilisateurs
\item Encadrement d'équipes pluridisciplinaires (design, ingénierie, marketing) en méthode Agile/Scrum
\item Forte expertise en discovery, tests utilisateurs, OKR et analytics produit
\item Passionnée par la création de solutions simples et impactantes pour résoudre des problèmes complexes
\end{itemize}}
\section{Formation}

\cventry{sept. 2012–juin 2015}
        {Master, Management de l'Innovation — Spécialisation Product Management, Score : 3.8}
        {HEC Paris}
        {\url{https://www.hec.edu}}
        {}
        {\begin{itemize}
\item Développement d'une double compétence stratégie et digital
\item Projet de fin d'études : création d'un prototype d'application mobile remarqué par un comité d'entrepreneurs
\item Échange universitaire à l'Université de Montréal (HEC Montréal) orienté vers le commerce électronique
\end{itemize}
\textbf{Cours} : Stratégie marketing et innovation, Gestion de projet et méthodes agiles, Analyse de données pour la décision, Expérience client et design thinking, Entrepreneuriat et business model}
\section{Expérience professionnelle}

\cventry{janv. 2022–Present}
        {Lead Product Manager}
        {Alto}
        {\url{https://www.alto.io}}
        {}
        {\begin{itemize}
\item Définition et exécution de la feuille de route produit d'une plateforme SaaS B2B dédiée aux paiements internationaux
\item Animation d'une équipe produit de 4 PM et coordination avec 6 équipes de développement
\item Mise en place d'un cadre de priorisation basé sur l'impact business et la valeur client, réduisant les délais de mise sur le marché de 30 \%
\item Conduite de plus de 50 entretiens utilisateurs et tests d'utilisabilité par trimestre pour éclairer les décisions produit
\item Pilotage des OKR produit et présentation des résultats à la direction générale et au conseil d'administration
\end{itemize}
\textbf{Mots-clés} : SaaS, Fintech, Leadership, Feuille de route, OKR, Data driven}

\cventry{janv. 2020–déc. 2021}
        {Product Manager}
        {Numera}
        {\url{https://www.numera.fr}}
        {}
        {\begin{itemize}
\item Gestion du cycle de vie complet de l'application mobile de gestion de finances personnelles (2,5 M d'utilisateurs)
\item Collaboration étroite avec le design et l'ingénierie pour livrer des fonctionnalités majeures : agrégation bancaire, scores financiers, notifications prédictives
\item Optimisation du parcours d'onboarding, avec une augmentation de +18 \% de l'activation à 7 jours
\item Expérimentations A/B en continu et analyse des entonnoirs de conversion via Amplitude et SQL
\end{itemize}
\textbf{Mots-clés} : Mobile, Tests A/B, Amplitude, SQL, Recherche utilisateur}

\cventry{mars 2018–déc. 2019}
        {Product Owner}
        {Lifya}
        {\url{https://www.lifya.com}}
        {}
        {\begin{itemize}
\item Conduite de la refonte du parcours d'achat du site e-commerce Lifya (CA annuel de 40 M€)
\item Rédaction des user stories, priorisation du backlog et animation des cérémonies Agile pour 3 équipes de développement
\item Coordination avec les équipes marketing, logistique et support pour améliorer la satisfaction client (NPS +24 points)
\item Mise en place des premiers indicateurs produit et des tableaux de bord de suivi de performance
\end{itemize}
\textbf{Mots-clés} : E-commerce, Scrum, Backlog, NPS, KPIs}
\section{Langues}

\cvline{Français}{Niveau natif ou bilingue \hfill \textbf{Mots-clés} : Langue maternelle}
\cvline{Anglais}{Niveau professionnel complet \hfill \textbf{Mots-clés} : TOEIC 985, Travail quotidien en anglais}
\cvline{Espagnol}{Niveau de travail limité \hfill \textbf{Mots-clés} : Niveau opérationnel}
\section{Compétences}

\cvline{Gestion de produit}{Expert \hfill \textbf{Mots-clés} : Découverte produit, Feuille de route, Priorisation, OKR, Spécifications fonctionnelles}
\cvline{Méthodes agiles}{Expert \hfill \textbf{Mots-clés} : Scrum, Kanban, Lean Startup, Agilité à l'échelle}
\cvline{Analyse de données}{Avancé \hfill \textbf{Mots-clés} : SQL, Amplitude, Mixpanel, Tableau de bord, Tests A/B}
\cvline{Design et recherche utilisateur}{Avancé \hfill \textbf{Mots-clés} : Recherche utilisateur, Entretiens terrain, Personas, Figma, Tests utilisateurs}
\section{Récompenses}

\cventry{nov. 2023}
        {Grand Prix Innovation - Bpifrance}
        {Prix du meilleur produit SaaS}
        {}
        {}
        {Récompense attribuée à la refonte de la plateforme Alto, saluée pour son expérience utilisateur et son impact sur la satisfaction client.}
\section{Certificats}

\cventry{mars 2021}
        {Pragmatic Institute}
        {Certificat Product Management - Niveau avancé}
        {\url{https://www.pragmaticinstitute.com}}
        {}
        {}

\cventry{août 2017}
        {Scrum Alliance}
        {Certified Scrum Product Owner (CSPO)}
        {\url{https://www.scrumalliance.org/certifications/cspo}}
        {}
        {}
\section{Publications}

\cventry{juin 2023}
        {Priorisation produit : méthodes RICE et WSJF comparées}
        {Le Product Blog}
        {\url{https://www.productblog.fr/priorisation-rice-wsjf}}
        {}
        {Article de synthèse sur les méthodes de priorisation produit, illustré par des cas pratiques issus du secteur fintech. Partage des critères d'impact, d'effort et de confiance pour aider les équipes à arbitrer efficacement.}
\section{Références}

\cventry{\emaillink[claire.dubois@alto.io]{claire.dubois@alto.io}}
        {Directrice produit chez Alto}
        {Claire Dubois}
        {+33 1 45 67 89 10}
        {}
        {Camille est une product manager exceptionnelle : elle sait combiner vision stratégique, rigueur d'exécution et intelligence relationnelle. Elle a considérablement élevé le niveau de notre organisation produit.}
\section{Projets}

\cventry{janv. 2023–déc. 2023}
        {Application personnelle de gestion budgétaire utilisée par plus de 500 personnes.}
        {BudgetFlow}
        {\url{https://www.budgetflow.fr}}
        {}
        {\begin{itemize}
\item Conception et développement d'un MVP de suivi budgétaire en no-code (Bubble, Airtable)
\item Réalisation d'entretiens utilisateurs et de tests d'empathie pour définir les fonctionnalités clés
\item Déploiement d'une première version publique et itérations basées sur les retours de la communauté
\end{itemize}
\textbf{Mots-clés} : No-code, MVP, Bubble, Airtable, Retours utilisateurs}
\section{Centres d'intérêt}

\cvline{Sport}{Course à pied, Yoga, Randonnée}
\cvline{Technologie}{Intelligence artificielle, No-code, Veille technologique}
\cvline{Lecture}{Essais économiques, Science-fiction}
\section{Bénévolat}

\cventry{janv. 2021–déc. 2023}
        {Mentor produit}
        {Paris Product Club}
        {\url{https://www.parisproductclub.fr}}
        {}
        {\begin{itemize}
\item Accompagnement de startups en phase d'amorçage sur leur stratégie produit
\item Animation d'ateliers de découverte et de priorisation pour des équipes fondatrices
\item Participation à des sessions de coaching collectif sur les méthodes Lean et Agile
\end{itemize}}

\end{document}